\documentclass[12pt]{article}

% ── Packages ──────────────────────────────────────────────────────────────
\usepackage[margin=1in]{geometry}
\usepackage{amsmath, amssymb}
\usepackage{graphicx}
\usepackage{booktabs}
\usepackage{natbib}
\usepackage{hyperref}
\usepackage{xcolor}
\usepackage{soul}          % \hl{} for TODO highlights
\usepackage{enumitem}
\usepackage{array}
\usepackage{caption}
\usepackage{subcaption}
\usepackage{multirow}
\usepackage{lineno}        % line numbers for review
\linenumbers

% ── TODO command ──────────────────────────────────────────────────────────
\newcommand{\todo}[1]{\hl{\textbf{[TODO: #1]}}}
\newcommand{\note}[1]{{\color{blue}\textit{[Note: #1]}}}

% ── Title block ───────────────────────────────────────────────────────────
\title{\textbf{From Regulatory Text to Shortage Risk:\\
  An NLP and LLM Framework for Pharmaceutical Quality Risk Prediction}}

\author{
  Amirreza Sahebi-Fakhrabad\thanks{Corresponding author.
    Poole College of Management, North Carolina State University.
    Email: \texttt{asahebi@ncsu.edu}} \and
  Shailesh Divey\thanks{Fisher College of Business, The Ohio State University.} \and
  John V.\ Gray\footnotemark[2] \and
  Robert Handfield\footnotemark[1] \and
  Amir Sadeghi\footnotemark[1]
}

\date{\today \\ \vspace{4pt}
  \textit{Preliminary draft — prepared for INFORMS Healthcare 2026.
  Please do not cite or circulate without permission.}}

% ══════════════════════════════════════════════════════════════════════════
\begin{document}
\maketitle

% ── Abstract ──────────────────────────────────────────────────────────────
\begin{abstract}
Drug shortages driven by manufacturing quality failures impose severe costs on patients,
hospitals, and supply chains.
Yet the rich unstructured regulatory text that FDA generates---inspection observations,
warning letters, recall narratives, and import refusal records---remains largely
untapped as a predictive resource.
This paper presents a multi-source regulatory intelligence framework integrating five
FDA databases covering 129 pharmaceutical manufacturing facilities across 14
high-priority generic Active Pharmaceutical Ingredients (APIs).
We construct a facility-level feature matrix combining structured regulatory outcomes
with NLP-derived text signals, enabling systematic prediction of downstream supply
disruptions.
Using NLP on structured regulatory text and an LLM extraction pipeline
(GPT-5-mini with strict JSON schema output) for Form 483 PDF processing, we
generate interpretable features capturing violation types, graded severity,
failure scope, escalation patterns, and latent firm behavior across two
theoretical frames: \emph{Agency Theory} (negligent process failures vs.\
persistent organizational misconduct) and \emph{Remediation Economics}
(capital-intensive vs.\ cultural root causes).
Preliminary results reveal actionable early-warning signals: specific regulatory
citation patterns strongly predict the most severe inspection outcomes
(Pearson $r = 0.51$); contamination and manufacturing practice failures account for
94\% of drug recalls in our sample; LLM severity grading separates the most severe
inspection outcomes (Critical+Major share 20 percentage points higher in
Official-Action-Indicated inspections than Voluntary-Action-Indicated); and
16\% of observations repeat citation language from an earlier inspection of the
same facility, identifying chronic violators.
\todo{Update shortage-model importance claims after re-running the prediction
pipeline with the revised raw text features.}
These signals can be detected months to years before shortages materialize, enabling
proactive supply chain risk management.
\end{abstract}

\textbf{Keywords:} drug shortage prediction, FDA Form 483, manufacturing quality,
large language model, NLP, pharmaceutical supply chain, quality risk index

\newpage

% ══════════════════════════════════════════════════════════════════════════
\section{Introduction}
\label{sec:intro}

Generic drugs account for roughly 90\% of U.S. prescriptions filled
\citep{fda2023generic}, yet generic drug shortages are an enduring public
health problem that disrupts patient care, inflates costs, and strains
hospital operations \citep{fox2014shortage}.
Shortages arise from an interplay of demand concentration, supply chain
fragility, and latent manufacturing quality failures that may not surface
until facilities fail FDA inspections or recall products.

The FDA's regulatory data—inspection outcomes (OAI/VAI/NAI classifications),
warning letters, recalls, and import refusals—provide structured signals of
quality risk, but they are largely lagging indicators that reflect failures
which have already materialized.
FDA Form 483 letters are issued at the \emph{conclusion} of an inspection and
document specific deficiencies before they escalate to warning letters or recalls.
Form 483 observations thus constitute a \emph{leading} signal of manufacturing
quality degradation, but their unstructured narrative format has prevented
systematic quantitative use.
Existing structured databases treat all compliance failures as binary flags,
masking the underlying organizational behavior.

Two theoretical frames guide our analysis.
\emph{Agency Theory}: we operationalize latent firm behavior by using regulatory
text to distinguish violations consistent with negligent process failures from
those indicative of persistent organizational misconduct.
\emph{Remediation Economics}: by categorizing root causes as capital-intensive
(equipment, infrastructure) versus cultural (procedures, data integrity), we
generate testable predictions on recovery dynamics---with cultural violations
expected to produce more persistent and volatile disruption patterns.

Prior work has established that FDA inspection outcomes predict downstream drug
shortages \citep{Stomberg2018, Wang2025MSOM} and that machine learning can
forecast shortages from pharmacy demand signals \citep{Liu2021, Pall2023}.
However, the rich narrative text that inspectors generate in Form 483
observations has not been used as a predictive input, and no study has directly
linked inspection-text signals to facility-level patient harm measured through
adverse event reporting \citep{Pazhayattil2020}.

This paper makes five contributions:
\begin{enumerate}[noitemsep]
  \item A novel multi-source regulatory text dataset covering 1,280 regulatory
        events across 129 facilities and five FDA databases, linked by FEI.
  \item A validated LLM extraction pipeline (Claude with strict JSON schema
        output, validated on a stratified sample before full deployment) that
        converts Form 483 observation narratives into structured, reproducible
        risk signals---with LLMs substantially outperforming rule-based
        extraction on contamination and patient-risk inference
        \citep{Li2026LLMDevice}.
  \item A time-aware facility-level feature panel: cumulative text-signal
        snapshots at each (FEI, inspection date), preventing future-information
        leakage in prediction, plus a cross-inspection repeat detector that
        identifies facilities cited for the same deficiency across inspections
        years apart.
  \item Empirical validation that LLM-graded severity tracks real inspection
        outcomes (Critical+Major share separates OAI from VAI inspections by
        20 percentage points), and that text-derived features carry information
        absent from structured databases.
  \item Discovery of two distinct failure modes encoded in 483 text: governance
        failures (quality system violations) predict regulatory escalation but
        not patient harm; technical failures (lab controls, data integrity)
        predict elevated serious adverse event rates but evade regulatory
        escalation---a distinction invisible to structured inspection outcome
        data alone.
\end{enumerate}

% ══════════════════════════════════════════════════════════════════════════
\section{Related Work}
\label{sec:lit}

\subsection{NLP and LLMs for FDA Regulatory Text}

Applying machine learning to FDA regulatory documents has concentrated on drug
labels and adverse-event reports.
Fine-tuned BERT models (PubMedBERT, BioBERT) reached F1 $\approx 0.90$ on
adverse drug reaction extraction from drug labels \citep{Tatonetti2025OnSIDES,
Bayer2021} and 0.84--0.87 MCC on drug-induced liver injury classification
\citep{Wu2021BERT}.
This domain-adaptation advantage has narrowed sharply: \citet{Li2026LLMDevice}
demonstrated that GPT-4-class models extract structured data from FDA medical
device adverse-event reports at 93--97\% human agreement \emph{without
fine-tuning}, and \citet{Hassani2025} found GPT-4o achieves 89\% precision on
regulatory text classification without task-specific training.
Applied to FDA recall narratives, multi-task BERT achieves 0.963 accuracy on
recall triage \citep{Atalay2026RecallRisk}, confirming that FDA regulatory
free text supports automated classification at scale.

On FDA inspection text specifically, \citet{Pazhayattil2020} conducted the
only published quantitative study of Form 483 data using pre-structured
citation frequency counts (2014--2018)---not natural language processing of
observation narratives.
Warning-letter text has been analyzed descriptively for domain and data
integrity trends \citep{Rathore2022, Park2025DataIntegrity, Kwiecinski2024}
but not used for feature extraction or prediction.
No prior study applies LLM extraction to Form 483 observation narratives to
generate structured quality-risk features.

\subsection{Inspection Outcomes and Drug Quality Events}

\citet{Stomberg2018} provided the first regression-based evidence that
inspection and citation intensity are significantly associated with drug
shortage rates.
\citet{Wang2025MSOM} extended this with instrumental variable analysis: OAI
outcomes predict a 96.4\% lower shortage likelihood in the subsequent
12~months.
Interpreted together, inspections appear to trigger corrective action that
stabilizes supply---but the \emph{severity of the underlying quality problem
before inspection} is not captured by binary OAI/VAI classifications.
Machine learning shortage prediction achieves C-statistic 0.69--0.93 using
pharmacy demand signals and drug characteristics \citep{Liu2021, Pall2023};
neither model incorporates inspection text or quality features.
\citet{Kosmas2023} modeled FDA inspection timing as a POMDP to maximize
value of inspecting high-risk facilities; their model assumes quality is
revealed by inspection but cannot observe it from text.
We provide the supply-side text features these models assume but cannot
currently observe.

\subsection{Adverse Events as a Manufacturing Quality Outcome}

FAERS adverse event reports are widely used for disproportionality analysis
but rarely deployed as a \emph{facility-level manufacturing quality metric}.
\citet{Sardella2021} provided the conceptual foundation: the nitrosamine
contamination crisis and heparin adulteration scandal were both first detected
through pharmacovigilance signal clustering tracing to facility-level
manufacturing decisions.
\citet{Brown2020} called for systematic use of FAERS for facility-level
quality surveillance and noted that the FDA already folds FAERS ``hazard
signals'' into its risk-based facility site-selection model for inspections.
\citet{Rahman2017FAERS} and \citet{Alatawi2018} developed authorized-generic
comparison methodology to control for perception bias in manufacturer-level
FAERS comparisons; our design largely sidesteps this confound because we
compare facilities manufacturing the \emph{same} API.
Consistent with \citet{Potter2025}, we treat FAERS counts as a predictive
correlate of quality, not a causal indicator.

We validate FAERS as a quality proxy against Valisure independent chemical
testing: drugs with higher Valisure analytical failure rates show
significantly more serious FAERS adverse events per million Medicaid
prescriptions (Spearman $\rho = 0.86$, $p = 0.014$, $n = 7$ drugs),
supporting \citet{Brown2020} and \citet{Sardella2021} empirically.
No prior study links FDA inspection text signals to facility-level FAERS
counts, or validates FAERS as a quality proxy against third-party chemical
testing.

% ══════════════════════════════════════════════════════════════════════════
\section{Data}
\label{sec:data}

We construct a unified regulatory event timeline for 129 facilities supplying
14 generic APIs, linked via FDA Facility Establishment Identifier (FEI) across
five sources yielding 1,280 total regulatory events (Table~\ref{tab:data}).

\begin{table}[h!]
\centering
\caption{Data sources and regulatory event coverage}
\label{tab:data}
\begin{tabular}{p{3.8cm} p{4.2cm} p{4.4cm}}
\toprule
\textbf{Source} & \textbf{Events / Coverage} & \textbf{Key Text Field} \\
\midrule
FDA Inspections + Citations & 1,019 inspections; 961 citations (127/129 FEIs) & CFR Short \& Long Description \\
FDA Form 483 PDFs & 38 FEIs with PDFs; 622 observations & Full observation narrative (unstructured) \\
FDA Warning Letters & 17 letters; 14 FEIs & Violation narrative; close-out status \\
FDA Drug Recalls & 49 API-matched; 18 FEIs & Reason for Recall (free text) \\
OASIS Import Refusals & 117 API-matched; 22 FEIs & Refusal charge codes; lab analysis flags \\
\bottomrule
\end{tabular}
\end{table}

Quality and shortage ground truth is drawn from three independent sources:
(1)~Valisure independent analytical testing scores for 14 APIs (2020--2024);
(2)~Merative MarketScan 2024 Commercial and Medicare pharmacy claims, which
captures dispensing volume anomalies and therapeutic switching as real-world
shortage signals; and
(3)~the University of Utah Drug Information Service (UUDIS) shortage database,
a clinically curated registry of active and historical shortage events.

\subsection{FDA Form 483 Observations}

FDA Form 483 is issued to facility management at the close of an inspection
when an investigator observes conditions that may violate the Food, Drug, and
Cosmetic Act (21 CFR Parts 210--211).
Each letter contains numbered observations written in investigator narrative
prose; they do not follow a templated structure beyond the opening
CFR-paraphrase sentence, making them the richest---but least structured---source
of regulatory intelligence.

We obtained 483 PDFs from FDA's public disclosure portal for 38 of our 129
facilities.
PDFs were converted to text via OCR (Tesseract 200 DPI) and parsed into
622 individual observations using regex splitting on ``OBSERVATION $N$'' headers.

\subsection{Sample}

The 14 APIs are: Ampicillin, Ampicillin/Sulbactam, Atorvastatin, Bupropion,
Calcium Gluconate, Lisinopril, Magnesium Sulfate, Metformin, Metoprolol,
Metronidazole, Pantoprazole, Potassium Chloride, Tacrolimus, and Vancomycin.
FEI-to-drug mapping uses the Valisure API--FEI crosswalk (193 unique FEI--drug
pairs; 129 unique FEIs).
The predictive panel covers 14 drugs $\times$ 10 years (2015--2024),
yielding 140 drug--year observations with 21 shortage-onset events (15.1\%).

% ══════════════════════════════════════════════════════════════════════════
\section{Methodology}
\label{sec:methods}

We apply three layers of text analysis, progressing from rule-based NLP to
LLM extraction.

\paragraph{Layer 1: NLP on Structured Regulatory Text.}
CFR citation Short Descriptions (230 unique values across 961 citations) are
aggregated into six regulatory domains---Laboratory/QC Controls, Production
Controls, Sterile Manufacturing, Documentation/Data Integrity,
Facility/Equipment, and Materials/Components---via CFR section mapping.
Recall reason free text is classified into nine categories (contamination,
cGMP failure, wrong potency, data integrity, etc.) using regex pattern matching.
Warning letter text is coded for violation domain and close-out resolution
status, providing a severity-escalation signal beyond routine citations.

\paragraph{Layer 2: LLM Pipeline for Form 483 PDF Extraction.}
Form 483 documents contain the richest regulatory intelligence---specific
measurements, equipment identifiers, and investigator narratives---but exist
only as unstructured PDFs (38 of 82 are scanned images requiring OCR).
The pipeline is: OCR (pytesseract) $\to$ observation-boundary chunking on
``OBSERVATION $N$'' headers $\to$ GPT-5-mini extraction under a strict JSON
schema $\to$ field validation with categorical coercion and an evidence guard.
Extraction is idempotent (already-scored observations are skipped on re-run)
and the prompt was validated on a stratified 50-observation sample with four
automated checks before the full corpus run (Section~\ref{sec:methods}).

\paragraph{Layer 3: Quality Risk Index Construction.}
A per-facility feature matrix integrating all regulatory signals feeds
correlation analysis against Valisure quality scores and predictive modeling
of shortage onset.

\subsection{LLM Extraction Schema}

For each Form 483 observation we send the cleaned observation text to
\texttt{gpt-5-mini} (OpenAI) with a structured JSON schema prompt.
The prompt instructs the model to return exactly 15 fields:

\begin{itemize}[noitemsep]
  \item \textbf{Categorical}: violation category (8 values), severity tier
        (Critical/Major/Moderate/Minor, EU GMP deficiency-classification
        style), failure scope (SingleBatch/MultipleProducts/FacilityWide/%
        Unclear), root-cause type (Capital/Cultural/Mixed/Unclear),
        remediation signal (Strong/Partial/Weak/None).
  \item \textbf{Binary flags} (5): repeat finding, patient risk, data
        integrity failure, contamination/sterility risk, investigation failure.
  \item \textbf{Qualitative}: severity rationale, root-cause rationale,
        verbatim evidence quote, confidence score (0--1).
\end{itemize}

The severity tier is graded on \emph{documented product impact} rather than
the seriousness of the system failure: Critical requires that affected product
was released or distributed; Major requires an actual defect, failing result,
or unreliable data found on site; Moderate covers deficient procedures or
systems with no actual defect documented; Minor covers administrative gaps.
The patient-risk flag requires an explicit harm pathway (sterile/injectable
product with a sterility assurance failure, a confirmed defect in released
product, or release without required QA disposition); ``could affect quality''
does not qualify.
The evidence quote is subject to an \emph{evidence guard}: the returned quote
must appear verbatim in the source text; otherwise it is cleared.

\paragraph{Validation.}
Before the full corpus run, the prompt was validated on a stratified
50-observation sample (round-robin across all 38 facilities) with four
automated checks: (a) the severity distribution is graded rather than
top-heavy; (b) the Critical+Major share separates OAI from VAI inspections
(via Redica inspection classifications matched on FEI and inspection date);
(c) flag fire rates fall in an informative range; and (d) the scope
distribution is spread across classes.
All four checks passed (OAI--VAI separation $+$0.20).
An earlier prompt draft failed check (a)---84\% of sampled observations were
rated Major---and was revised to the documented-product-impact rubric above
before deployment.
\todo{Add independent human-rater agreement (Cohen's $\kappa$) on a
manually reviewed sample.}

\subsection{Facility-Level Aggregation}

We aggregate observation-level signals into \emph{time-aware cumulative
snapshots}: one row per (FEI, 483 inspection date), where each feature is
computed from all observations of that facility up to and including that
date.
The downstream shortage model joins on FEI and takes the most recent snapshot
before each prediction period, so no feature ever encodes information from a
future inspection.
Rather than hand-weighted composite indices, raw feature shares are passed
directly to the model, which learns weights from data: severity tier shares
(including the combined Critical$+$Major share), scope shares, violation
category shares (8 classes), root-cause and remediation shares, the five LLM
flag shares, and two regex flags with no LLM equivalent (OOS/OOT events and
prior-warning-letter references).

\paragraph{Cross-inspection repeat detection.}
The per-observation repeat flags only catch repeats the investigator
explicitly labels (``this is a repeat observation''), firing on 6.8\%
(regex) and 10.0\% (LLM) of observations.
We add a complementary signal that no per-observation method can produce:
each 483 observation opens with a standardized citation sentence (the
template text before ``Specifically,'').
When that sentence matches an observation from a strictly earlier inspection
of the same facility (token overlap $\geq 0.80$), the same deficiency was
cited again.
This flags 97 of 616 dated observations (15.7\%) as cross-inspection
repeats; the feature is time-safe by construction because it only looks
backward.
Full-text similarity fails at this task: shared GMP vocabulary gives
unrelated observations a $\sim$0.45 baseline token overlap, so matching is
restricted to the standardized citation sentence.

\subsection{Manufacturing Quality Risk Index (MQRI)}

The MQRI combines regulatory enforcement features and market quality signals:

\begin{equation}
  \text{MQRI}_i = \left( 0.70 \cdot D^{\text{reg}}_i
                         + 0.30 \cdot D^{\text{mkt}}_i \right) \times 100
\end{equation}

where $D^{\text{reg}}$ includes cumulative OAI count, OAI-predictive CFR
inspections, posted 483 citations, warning letters, inspection gap, import
refusals, and the contamination NLP flag; $D^{\text{mkt}}$ includes Class I
and all-class recall counts.
Feature weights are derived from Spearman $|\rho|$ with OAI rate across
$n = 127$ facilities.
In this paper we extend the MQRI by adding the 483 text severity signal
(Critical$+$Major share) as a regulatory domain feature with weight derived
from \todo{correlation with OAI rate or Valisure}.

\subsection{Shortage Prediction Model}

We construct a drug--year panel (14 APIs $\times$ 2015--2024; $n=140$ rows,
21 shortage-onset events) and fit grouped cross-validated benchmark models
(GroupKFold by drug, $k=5$) to predict $y_{t+1} = 1$ if shortage starts in
year $t+1$.
Features at year $t$ include FAERS adverse-event signals, Redica
inspection/regulatory signals, Valisure quality scores (time-invariant), and
the LLM-derived raw text features (severity Critical$+$Major share, scope
shares, flag shares, cross-inspection repeat share), aggregated to
drug level by weighted mean across FEIs (weights: $n_{\text{obs scored}}$).
Recalls are excluded as predictors---timing analysis confirms they are
concurrent or lagging relative to shortage onset, not leading indicators.
We compare L2-penalized logistic regression and Random Forest (300 trees,
balanced class weights) with and without text features, reporting AUC and
the ablation delta.

% ══════════════════════════════════════════════════════════════════════════
\section{Results}
\label{sec:results}

\subsection{Descriptive Characterization of Form 483 Observations}

Table~\ref{tab:obs_desc} summarizes the LLM extraction results across
622 observations from 38 facilities.

\begin{table}[h!]
\centering
\caption{Observation-level signal distribution ($n = 622$ observations,
         38 facilities)}
\label{tab:obs_desc}
\begin{tabular}{lrr}
\toprule
\textbf{Signal} & \textbf{Count} & \textbf{Share (\%)} \\
\midrule
\multicolumn{3}{l}{\textit{Severity tier}} \\
\quad Critical  &  99 & 15.9 \\
\quad Major     & 188 & 30.2 \\
\quad Moderate  & 331 & 53.2 \\
\quad Minor     &   4 &  0.6 \\
\midrule
\multicolumn{3}{l}{\textit{Failure scope}} \\
\quad MultipleProducts & 374 & 60.1 \\
\quad FacilityWide     & 190 & 30.5 \\
\quad SingleBatch      &  46 &  7.4 \\
\quad Unclear          &  12 &  1.9 \\
\midrule
\multicolumn{3}{l}{\textit{Primary violation domain}} \\
\quad Production Controls  & 184 & 29.6 \\
\quad Lab Controls         & 159 & 25.6 \\
\quad Quality System       & 115 & 18.5 \\
\quad Buildings/Equipment  &  72 & 11.6 \\
\quad Records/Reports      &  53 &  8.5 \\
\quad Personnel            &  21 &  3.4 \\
\quad Packaging/Labeling   &  13 &  2.1 \\
\quad Other                &   5 &  0.8 \\
\midrule
\multicolumn{3}{l}{\textit{Root-cause type}} \\
\quad Cultural  & 271 & 43.6 \\
\quad Mixed     & 228 & 36.7 \\
\quad Capital   & 118 & 19.0 \\
\quad Unclear   &   5 &  0.8 \\
\midrule
\multicolumn{3}{l}{\textit{Remediation signal}} \\
\quad None    & 443 & 71.2 \\
\quad Partial & 103 & 16.6 \\
\quad Weak    &  74 & 11.9 \\
\quad Strong  &   2 &  0.3 \\
\bottomrule
\end{tabular}
\end{table}

\paragraph{Key finding 1: Severity grades on documented product impact.}
Under the documented-product-impact rubric, 15.9\% of observations are
Critical (defective product was released or distributed), 30.2\% Major
(an actual defect or failing result found on site), and 53.2\% Moderate
(deficient system, no defect documented).
The Critical share is striking: roughly one in six observations in this
sample documents defective product that reached the market.
The graded distribution---unlike a rubric where any serious-sounding GMP
failure rates highly---gives the tier discriminating power:
the Critical$+$Major share is 20 percentage points higher in inspections
that FDA classified OAI than VAI.

\paragraph{Key finding 2: Cultural root causes dominate.}
80.3\% of observations are classified as Cultural or Mixed root cause,
suggesting that the predominant drivers of CGMP violations in this sample
are management, training, and quality-culture failures rather than
capital/equipment gaps.
This has direct implications for corrective action: capital investment alone
is insufficient.

\paragraph{Key finding 3: Remediation is rarely mentioned.}
71.2\% of observations mention no corrective action whatsoever, and only
two observations (0.3\%) reference a strong corrective action.
\note{This is partly structural: 483 observations are investigator-written
findings; firm responses are submitted separately.
However, Partial/Weak mentions do appear when investigators note commitments
made on-site. The absence is itself informative.}

\subsection{Semantic Lift: LLM vs.\ Regex}

Table~\ref{tab:lift} compares LLM flag rates to deterministic regex flags
on the same observations.

\begin{table}[h!]
\centering
\caption{Semantic lift: LLM vs.\ regex flag rates ($n = 622$ observations)}
\label{tab:lift}
\begin{tabular}{lrrr}
\toprule
\textbf{Flag} & \textbf{Regex share} & \textbf{LLM share}
              & \textbf{LLM-only lift} \\
\midrule
Contamination  & 33.1\% & 58.5\% & \textbf{27.3\%} \\
Patient risk   & 19.9\% & 33.6\% & \textbf{21.2\%} \\
Data integrity &  5.3\% & 18.5\% & 13.8\% \\
Investigation  & 37.5\% & 40.0\% &  7.2\% \\
Repeat finding &  6.8\% & 10.0\% &  3.2\% \\
\bottomrule
\end{tabular}
\end{table}

The contamination and patient-risk flags show the largest semantic lift:
the LLM flags 27 and 21 percentage points of observations, respectively,
that keyword rules miss.
This is because these inferences require contextual reasoning
(e.g., a sterile product exposed to unclassified air $\Rightarrow$ patient
risk, even though the words ``patient'' and ``contamination'' never appear)
that regex cannot replicate.
The reverse error also occurs: for patient risk, 7.6\% of observations are
regex-only---the word ``patient'' appears but the LLM judged that no
explicit harm pathway exists under the strict prompt definition.

\subsection{Facility-Level Text Features}

\todo{Table of top/bottom 10 FEIs by severity\_critmajor\_share and
repeat\_cross\_insp\_share with firm name (or anonymized).
Add: correlation of severity\_critmajor\_share with n\_oai,
correlation with MQRI.}

\subsection{Validation: FAERS as a Manufacturing Quality Proxy}

Before linking 483 text signals to adverse event outcomes, we validate that
FAERS serious AE counts track manufacturing quality rather than pharmacological
effects alone.
For each of the seven generic drugs with published Valisure analytical quality
scores, we compute the fraction of ANDA-holders with score $< 90$ (Valisure
failure rate) and the drug's FAERS serious adverse event count per million
Medicaid prescriptions (using CMS SDUD prescription volume as denominator).
Spearman correlation across seven drugs yields $\rho = 0.86$, $p = 0.014$:
drugs with higher proportions of analytically deficient products generate more
serious adverse events per patient exposed, even after controlling for
prescribing volume (Table~\ref{tab:valisure_val}).
The SDUD-normalized metric is stronger than the raw AE count ($\rho = 0.79$),
indicating that prescription volume was attenuating, not inflating, the
quality signal.

\begin{table}[h!]
\centering
\small
\caption{Valisure failure rate vs.\ FAERS serious AEs per million Medicaid Rx
         ($n = 7$ drugs; Spearman $\rho = 0.86$, $p = 0.014$)}
\label{tab:valisure_val}
\begin{tabular}{lrr}
\toprule
\textbf{Drug} & \textbf{Valisure fail rate} & \textbf{Serious AEs / M Rx} \\
\midrule
Tacrolimus         & 69\% & 1{,}738 \\
Metformin          & 48\% &   133 \\
Metoprolol         & 46\% &    61 \\
Potassium chloride & 41\% &    21 \\
Magnesium sulfate  &  0\% &    19 \\
Lisinopril         & 17\% &    16 \\
Calcium gluconate  & 33\% &     7 \\
\bottomrule
\end{tabular}
\end{table}

\subsection{Two Failure Modes: Governance vs.\ Technical Failures}

We test whether the same text signals that predict regulatory escalation also
predict patient harm, or whether inspection text encodes two distinct failure
pathways.
For each of 98 FEIs with 483 text features, we compute Spearman correlations
of each text signal against three outcomes: (1)~OAI inspection rate
(regulatory escalation), (2)~Class~I recall count (product-level regulatory
action), and (3)~mean FAERS serious AEs at lag $+$1~year (patient harm).

Results reveal a clear divergence (Table~\ref{tab:two_modes}):

\paragraph{Governance failures $\to$ regulatory action, not patient harm.}
\emph{Quality system violations} are the strongest predictor of OAI rate
across all ten text features ($\rho = +0.33$, $p < 0.001$) and show no
significant relationship with serious AEs ($\rho = -0.08$, $p > 0.4$).
Quality system observations cite inadequate oversight, missing review
procedures, and quality-unit authority failures---exactly the systemic
indicators FDA uses to justify OAI classification.
The regulatory response (mandatory CAPA, import alerts, product holds) limits
patient exposure, attenuating the harm signal.
\emph{Repeat observations} trend toward both OAI and recalls ($\rho \approx
+0.20$ for both), consistent with repeat citations building a regulatory
history that accelerates escalation.

\paragraph{Technical failures $\to$ patient harm, evading regulatory escalation.}
\emph{Lab controls violations}---inadequate test methods, OOS investigation
failures, missing specifications---are the strongest predictor of FAERS
serious AEs ($\rho = +0.32$, $p < 0.01$) and show essentially no relationship
with OAI rate ($\rho = +0.01$).
Lab controls are the gatekeeping function: when they fail, defective product
passes QC without generating the governance-level signals that trigger OAI.
\emph{Data integrity violations} follow the same pattern ($\rho = +0.22$ for
AEs, $p < 0.05$; $\rho = +0.10$ for OAI, $p > 0.3$).
DI failures are by design difficult to detect and slow to escalate---falsified
records suppress the evidence trail that inspectors use to assign OAI.

\paragraph{Removed products.}
\emph{Contamination flags} show negative correlations with both OAI
($\rho = -0.05$) and AEs ($\rho = -0.19$), consistent with contamination
findings triggering recalls that remove product from the market before
widespread patient harm materializes.

\begin{table}[h!]
\centering
\small
\caption{Spearman $\rho$ — each 483 text signal vs.\ regulatory action and
         patient harm outcomes ($n = 98$ FEIs)}
\label{tab:two_modes}
\begin{tabular}{lrrr}
\toprule
\textbf{Text signal} & \textbf{OAI rate} & \textbf{Class~I recalls}
                     & \textbf{Serious AEs (lag+1)} \\
\midrule
\multicolumn{4}{l}{\textit{Governance failures → regulatory action}} \\
\quad Quality system violations & $+0.33$*** & $-0.02$ & $-0.08$ \\
\quad Repeat observations       & $+0.20$    & $+0.19$ & $+0.13$ \\
\quad Investigation failures    & $+0.17$    & $+0.04$ & $+0.13$ \\
\midrule
\multicolumn{4}{l}{\textit{Technical failures → patient harm}} \\
\quad Lab controls violations   & $+0.01$    & $+0.18$ & $+0.32$** \\
\quad Data integrity violations & $+0.10$    & $-0.09$ & $+0.22$* \\
\quad Scope: facility-wide      & $+0.10$    & $-0.10$ & $+0.10$ \\
\midrule
\multicolumn{4}{l}{\textit{Removed products}} \\
\quad Contamination flag        & $-0.05$    & $-0.07$ & $-0.19$ \\
\quad Patient risk flag         & $+0.14$    & $-0.03$ & $-0.06$ \\
\bottomrule
\multicolumn{4}{l}{\footnotesize * $p<0.05$; ** $p<0.01$; *** $p<0.001$} \\
\end{tabular}
\end{table}

\subsection{Predictive Validity: Text Features vs.\ Valisure Outcomes}

\todo{Run correlation: severity\_critmajor\_share vs.\ Valisure dissolution
difference factor (2024).
Run correlation: contamination\_llm\_share vs.\ Valisure NDMA/DMF.
Compare to MQRI-only baseline.}

\subsection{Predictive Validity: LLM Text Features in Shortage Prediction}

\note{The benchmark below was estimated with the previous feature set
(composite text indices, 3-tier severity). The extraction pipeline has since
been revised: 4-tier severity, scope field, 5 binary flags, cross-inspection
repeat detection, and raw feature shares in place of the four composite
indices. The model must be re-run before these numbers are cited.}

\todo{Re-run m07/m09 with the revised text features
(severity\_critmajor\_share, scope shares, flag shares,
repeat\_cross\_insp\_share) and replace the table, figure, and importance
claims below.}

Table~\ref{tab:model} reports grouped cross-validated AUC for the full model
and the ablation without text features.

\begin{table}[h!]
\centering
\caption{Shortage prediction benchmark — GroupKFold CV by drug
         ($n=126$ usable rows; 19 events; 14 drugs, 2015--2024).
         \todo{Stale — re-run with revised text features.}}
\label{tab:model}
\begin{tabular}{lccc}
\toprule
\textbf{Model} & \textbf{AUC (with text)} & \textbf{AUC (without text)} & \textbf{$\Delta$AUC} \\
\midrule
L2 Logistic Regression & 0.545 & 0.610 & $-$0.065 \\
Random Forest          & 0.590 & 0.538 & $+$0.052 \\
\bottomrule
\end{tabular}
\end{table}

Random Forest AUC improves by $+$0.052 with text features.
The logistic regression result is negative ($-$0.065), consistent with
high-dimensional collinearity at $n = 19$ events---logit is not robust
to correlated predictors in small samples.

\begin{figure}[h!]
\centering
\includegraphics[width=0.75\textwidth]{../../Data/99 - Outputs - Shortage Prediction/outputs/figures/feature_importance_valisure_clean.png}
\caption{Random Forest feature importance for predicting next-year drug shortage (14 APIs, 2015--2024, GroupKFold CV by drug). Teal bars = LLM-extracted 483 text features; navy bars = structured data features. \todo{Regenerate with revised raw text features.}}
\label{fig:fi}
\end{figure}

% ══════════════════════════════════════════════════════════════════════════
\section{Discussion}
\label{sec:discussion}

\paragraph{Form 483 text as a leading quality signal.}
LLM classification of 483 observation narratives produces signals that
structured compliance databases do not capture: graded severity that tracks
FDA inspection outcomes (Critical$+$Major share 20pp higher in OAI than VAI
inspections), failure scope, remediation commitment, and chronic-violator
identification via cross-inspection repeat detection (15.7\% of
observations repeat citation language from an earlier inspection of the
same facility).
\todo{Restate the shortage-model importance finding after the re-run with
raw text features.}
The semantic lift analysis shows that LLMs detect contamination-control risk
and patient-risk implications in 27 and 21 percentage points more
observations, respectively, than keyword rules---confirming that contextual
reasoning, not pattern matching, is required to extract these signals.

\paragraph{Root cause composition and Remediation Economics.}
80.3\% of observations carry Cultural or Mixed root causes, consistent with
Agency Theory's prediction that persistent non-compliance reflects management
behavior rather than capital gaps.
CAPA programs that target equipment or procedures may therefore be
systematically underspecified; cultural-root-cause facilities are predicted to
show heavier-tailed recovery dynamics.
\todo{Add survival analysis results when available.}

\paragraph{Research investigations for full paper.}
Three empirical investigations structure the full analysis:
(1)~\emph{Preemptive market dynamics}: do compliance signals trigger early
dispensing drops, price spikes, or therapeutic switching before official
shortages are declared (MarketScan outcome)?
(2)~\emph{Endogeneity of recovery}: do capital violations yield predictable
resolution windows while cultural violations produce relapsing shortage durations
(survival analysis)?
(3)~\emph{Supply network resilience}: do disruptions at high-risk facilities
propagate through the API supply network or are they absorbed by competing
suppliers?

\paragraph{Limitations.}
\begin{itemize}[noitemsep]
  \item \textit{Coverage}: 483 PDFs available for only 38 of 129 FEIs, limiting
        text-feature reach; structured citation signals cover 127/129.
        Drug-level text indices therefore reflect only the subset of FEIs with
        publicly available inspection documents — for example, three of four
        Ampicillin-manufacturing FEIs have no 483 PDFs on file, so the
        drug-level average is derived from a single facility.
  \item \textit{Cross-inspection repeat detection}: the template-sentence
        matcher catches deficiencies re-cited under the same standardized
        citation language (e.g., FEI 3011905047, a sterile injectable
        manufacturer, was cited for the identical employee-training
        deficiency in 2019 and again in 2024).
        It will miss repeats that an investigator cites under a different
        CFR paraphrase, so the 15.7\% cross-inspection repeat rate is a
        lower bound.
  \item \textit{Sample size}: 14 drugs and 19 shortage events yield wide
        confidence intervals; all findings are exploratory.
  \item \textit{LLM reliability}: manual review of a sample confirms strong
        precision on violation category and severity; contamination flags
        occasionally over-trigger on training-context observations.
  \item \textit{Selection}: Valisure tests a non-random subset of generics.
  \item \textit{Causality}: we estimate predictive associations, not causal effects.
\end{itemize}

% ══════════════════════════════════════════════════════════════════════════
\section{Conclusion}
\label{sec:conclusion}

We present a multi-source regulatory intelligence framework that integrates
five FDA databases and an LLM extraction pipeline to convert Form 483
narrative observations into structured, predictively informative quality-risk
features.
Key findings: (1) LLM-graded severity tracks real inspection outcomes
(Critical$+$Major share 20pp higher in OAI than VAI inspections), and 15.9\%
of observations document defective product that reached the market;
(2) 15.7\% of observations repeat citation language from an earlier
inspection of the same facility, identifying chronic violators that
per-observation flags miss; (3) 80.3\% of violations carry Cultural or
Mixed root causes, with direct implications for CAPA program design;
(4) contamination and cGMP failures account for 94\% of API-matched recalls.
\todo{Add shortage-model finding after the re-run with raw text features.}
The pipeline is fully automated and reproducible from publicly available FDA
data, offering a scalable framework for real-time manufacturing quality
surveillance and proactive supply chain risk management---with direct
applications in hospital formulary management, GPO contract decisions, and
FDA surveillance prioritization.
To our knowledge, this is the first study to link FDA regulatory inspection
text to pharmacy claims data as a validation of supply disruption, closing
the loop from upstream manufacturing quality signals to patient-level drug
access.

% ══════════════════════════════════════════════════════════════════════════
\bibliographystyle{agsm}
\bibliography{references}

% ══════════════════════════════════════════════════════════════════════════
\appendix

\section{LLM Flag Definitions}
\label{app:flags}

\paragraph{Extraction model.} \texttt{gpt-5-mini} with structured JSON output
schema (strict mode). Max output tokens = 4000. Rate-limited requests are
retried with linear back-off; failed rows are re-scored on the next run
rather than silently imputed.

\paragraph{Evidence guard.} The \texttt{evidence\_quote} field must appear
verbatim in the source observation text. If the model paraphrases, the field
is cleared to empty string rather than accepting an inaccurate citation.

\paragraph{Flag definitions} are as specified in the extraction prompt
(see \texttt{01\_extract\_observation\_signals.py}):

\begin{description}[leftmargin=2em]
  \item[\texttt{repeat\_flag\_llm}] True only when the observation explicitly
    states this is a repeat or previously-cited finding from a prior inspection.
    Not true for repeated examples within the same current observation.
  \item[\texttt{patient\_risk\_flag\_llm}] True only when an explicit harm
    pathway exists: (a) sterile/injectable product with a contamination or
    sterility assurance failure, (b) a confirmed quality defect in released
    or distributed product, or (c) product released without required QA
    disposition. ``Could affect quality'' does not qualify.
  \item[\texttt{data\_integrity\_flag\_llm}] True only for explicit data
    trustworthiness failures (falsification, backdating, audit-trail problems).
    Not true for ordinary incomplete documentation.
  \item[\texttt{contamination\_flag\_llm}] True for actual contamination or
    clear contamination-control risk including sterility assurance failures.
  \item[\texttt{investigation\_flag\_llm}] True only for an explicit failed,
    missing, or inadequate investigation of a concrete event.
\end{description}

\paragraph{Removed in prompt v2.} The former \texttt{systemic\_flag\_llm}
fired on 94\% of observations (uninformative) and was replaced by the
4-class \texttt{scope} field; \texttt{documentation\_flag\_llm} is covered
by the RecordsReports violation category.

\section{Severity Tier Definitions (Prompt v2)}
\label{app:severity}

The tier is decided by one question: what level of actual product impact
does the text document?

\begin{description}[leftmargin=2em]
  \item[Critical] Affected product was released or distributed (e.g.,
    contaminated lots distributed before the investigation was closed).
  \item[Major] An actual defect, failure, or unreliable result found at the
    facility but no evidence of release (e.g., particulate matter observed
    in lots; test results invalidated without quality approval).
  \item[Moderate] A deficient procedure, system, or practice with no actual
    product defect documented (e.g., media fills that do not mimic routine
    production; systems that allow data deletion without evidence it
    occurred). The default tier for most observations.
  \item[Minor] Documentation or administrative gap with no plausible product
    impact.
\end{description}

\end{document}
